% SGA 3，第 XI 篇；根据英文重建源翻译。
% 范围：命题 3.12、注 3.13 及完整直接证明；止于第 4 节之前。
% 权威来源：Polo--Gille Expo11.pdf，组件局部页 17--20；合卷页 739--742；印刷／原文页边页 154--156。
% Jacob Reinhold 的 CC BY 4.0 Markdown 谱系仅用于措辞比较；文字、公式、编号、注释及图均以 Polo--Gille PDF 为准。
% 精确续接／终点：第 4 节，Expo11 局部页 20／合卷页 742／印刷页 157。

\medskip\noindent\SGARefTarget{sga3:xi:proposition:3:12}\textbf{命题 3.12。} --- \textit{设 \(S\) 为预概形，\(G\) 为仿射 \(S\)-群预概形。}

\begin{enumerate}\item[\SGARefTarget{sga3:xi:proposition:3:12:item:a}\textit{(a)}]\textit{设 \[
  F:(\mathbf{Sch})^\circ/S\longrightarrow(\mathbf{Ens})
\] 为函子，使得对 \(S\) 上的每个 \(T\)，} \SGARefTarget{sga3:xi:display:3:12:finite-multiplicative-subgroups} \[
  F(T)=
  \left\{
  \begin{array}{c}
  \text{\(G_T\) 的乘法型子群}\\
  \text{在 \(T\) 上有限}
  \end{array}
  \right\}.
\]\textit{。假设 \(S\) 局部 Noether，或 \(G\) 在 \(S\) 上有限呈示。则 \(F\) 可表且在 \(S\) 上仿射。若 \(G\) 在 \(S\) 上有限呈示，则 \(F\) 在 \(S\) 上局部有限呈示。}

\item[\SGARefTarget{sga3:xi:proposition:3:12:item:b}\textit{(b)}]\textit{设 \(H\) 为有限乘法型 \(S\)-群预概形。则 \(\operatorname{Hom}_{S\text{-}\mathrm{gr}}(H,G)\) 可表且在 \(S\) 上仿射。若 \(G\) 在 \(S\) 上有限型（分别为有限呈示），则 \(\operatorname{Hom}_{S\text{-}\mathrm{gr}}(H,G)\) 亦然。}\end{enumerate}

\medskip\noindent\SGARefTarget{sga3:xi:remark:3:13}\textbf{注 3.13。} --- 除 \(\operatorname{Hom}_{S\text{-}\mathrm{gr}}\) 仿射的断言以及 \(G\) 在 \(S\) 上有限呈示的情形（此情形已足够我们使用）以外，\SGARefLink{sga3:xi:proposition:3:12}{3.12} 是 Hilbert 概形理论的直接推论（A.~Grothendieck，\textit{代数几何中的构造技术与存在定理：IV．Hilbert 概形〔Techniques de construction et théorèmes d'existence en Géométrie Algébrique: IV. Les Schémas de Hilbert〕}，Séminaire Bourbaki，1961 年 5 月，第~221 号）。事实上，只需 \(G\) 在 \(S\) 上拟射影。在情形 \textit{\SGARefLink{sga3:xi:proposition:3:12:item:a}{(a)}} 中，还可表示更大的函子 \[
  F_0(T)=
  \left\{
  \begin{array}{c}
  \text{\(G_T\) 的平坦、固有群子预概形，}\\
  \text{且在 \(T\) 上有限呈示}
  \end{array}
  \right\}.
\]。由 \SGARefLink{sga3:xi:proposition:3:1}{3.1} 的判据及 \SGARefLink{sga3:x:corollary:4-7}{第 X 篇，4.7}~\textit{(b)}，典范单态射 \(F\to F_0\) 为“开浸入”。因此，\(F_0\) 的可表性蕴含 \(F\) 可由 \(F_0\) 的开子预概形表示。在情形 \textit{\SGARefLink{sga3:xi:proposition:3:12:item:b}{(b)}} 中，只需假设 \(H\) 在 \(S\) 上射影且有限呈示。在两种情形中，得到的函子均在 \(S\) 上局部有限呈示。

在本篇中，\SGARefLink{sga3:xi:proposition:3:12}{3.12} 只是用于证明 \SGARefLink{sga3:xi:section:4}{下一节} 中一个关键结果的技术引理。因此我们将简述不使用 Hilbert 概形的一个简单直接证明。

先证明 \textit{\SGARefLink{sga3:xi:proposition:3:12:item:b}{(b)}}。只需不依赖于 \(H\) 和 \(G\) 上的群结构，证明 \(\operatorname{Hom}_S(H,G)\) 可表且具有对 \(\operatorname{Hom}_{S\text{-}\mathrm{gr}}(H,G)\) 所断言的附加性质。事实上，对 \(\operatorname{Hom}_S(H\times_S H,G)\) 应用同一结果，\(\operatorname{Hom}_S(H,G)\) 的子函子 \(\operatorname{Hom}_{S\text{-}\mathrm{gr}}(H,G)\) 可表示为涉及 \(G\)、\(\operatorname{Hom}_S(H,G)\) 和 \(\operatorname{Hom}_S(H\times_S H,G)\) 的有限逆极限，实际上是纤维积。此表达式留给读者。

写成 \[
  H=\operatorname{Spec}(B),
\]，其中 \(B\) 为 \(S\) 上的代数层，作为模层局部自由；这是论证对 \(H\) 所用的唯一假设。由于可表性问题在 \(S\) 上是局部的，可以假设 \(S\) 仿射，环为 \(A\)。还可以写成 \(G=\operatorname{Spec}(C)\)，其中 \(C\) 为 \(A\)-代数。

当 \(G=S[t]=G_0\)，其中 \(t\) 为不定元时，Hom 函子就是 \[
  T\longmapsto\Gamma(T,B_T).
\]。由于 \(B\) 局部自由，此函子由向量丛 \(V(B^\vee)\) 表示，其中 \(B^\vee\) 为对偶模层。当 \(G=S[(t_i)_{i\in I}]=G_0^I\)，其中不定元族不必有限时，可得\footnote{原文显示式印作 \(\operatorname{Hom}_S(H,G)^I=V(B^\vee)^I=V(B^{\vee(I)})\)。由于前一句令 \(G=G_0^I\)，此显示等式须明确写出因子 \(\operatorname{Hom}_S(H,G)=\operatorname{Hom}_S(H,G_0)^I\)。} \SGARefTarget{sga3:xi:display:3:13:hom-affine-space} \[
  \operatorname{Hom}_S(H,G)
  =\operatorname{Hom}_S(H,G_0)^I
  =V(B^\vee)^I
  =V\bigl(B^{\vee(I)}\bigr).
\]。

一般地，\(G\) 同构于形如 \(S[(t_i)_{i\in I}]\) 的概形的闭子概形；等价地，\(C\) 为某个 \(A\)-代数 \(A[(t_i)_{i\in I}]\) 的商。设 \((F_j)\) 为所取商的理想的一族生成元。将 \(S\) 用较小的仿射开集覆盖后，可以假设 \(B\) 为秩 \(n\) 的自由模。选取 \(B\) 的基 \((e_k)_{1\leq k\leq n}\)。写成 \[
  x_i=\sum_{k=1}^{n}x_{ik}e_k
\]，其中系数 \(x_{ik}\) 是某个未指定的 \(A\)-代数 \(A'\) 中的不定系数，将 \(F_j((x_i))\) 按此基展开。对每个 \(F_j\)，这给出 \(n\) 个以 \((x_{i,k})_{i\in I,\,1\leq k\leq n}\) 为变量、系数在 \(A\) 中的多项式 \[
  F_{j,k}\qquad(1\leq k\leq n)
\]。立即可知 \(\operatorname{Hom}_S(H,G)\) 由多项式环 \(A[(x_{i,k})]\) 对各 \(F_{j,k}\) 生成的理想取商所得环的谱表示。这就证明了 \textit{\SGARefLink{sga3:xi:proposition:3:12:item:b}{(b)}}。

现在证明 \textit{\SGARefLink{sga3:xi:proposition:3:12:item:a}{(a)}}。对每个通常意义下的有限交换群 \(M\)，设 \(F_M\) 为 \(F\) 的子函子，将所考虑的子群限制为 \(G_T\) 中的 \(M\) 型乘法型子群。类似于 \SGARefLink{sga3:xi:proposition:3:5}{3.5} 中的粘合论证（本应将问题再稍作分解后才陈述它）表明，只需证明每个 \(F_M\) 可表。于是 \(F\) 由各 \(F_M\) 的不交和表示，其中 \(M\) 遍历有限交换群的同构类。事实上，\(F\) 为层范畴中的相应和。

固定 \(M\)，以 \(F\) 代替 \(F_M\)。置 \(H=D_S(M)\)，设 \[
  F'=\operatorname{Imm}_{S\text{-}\mathrm{gr}}(H,G)
  \subset\operatorname{Hom}_{S\text{-}\mathrm{gr}}(H,G)
\] 为子函子，其在 \(T\) 处的值为所有作为闭浸入的 \(T\)-群同态 \(H_T\to G_T\) 的集合。由 \textit{\SGARefLink{sga3:xi:proposition:3:12:item:b}{(b)}}，\(\operatorname{Hom}_{S\text{-}\mathrm{gr}}(H,G)\) 由仿射 \(S\)-预概形表示。利用 \SGARefLink{sga3:ix:theorem:6:8}{第 IX 篇，6.8}，立即可知 \(F'\) 由其中的开闭子预概形表示；特别地，\(F'\) 在 \(S\) 上仿射。

考虑典范态射 \(F'\to F\)，它将每个单态射 \(H_T\to G_T\) 映为其像群。直接由定义，在 fpqc 层范畴中，\(F'\) 因而成为 \(F\) 上在 \[
  \Gamma_F=\Gamma_S\times_S F,
  \qquad
  \Gamma=\operatorname{Aut}_{\mathrm{gr}}(M).
\] 作用下的主齐性丛。由下降可知 \(F'\to F\) 可表：对每个态射 \(T\to F\)，其中 \(T\) 可表，预概形 \(T\times_F F'\) 在 \(T\) 上可表，实际上可由 \(\Gamma\) 作用下的主 Galois 丛表示。

由 \SGARefLink{sga3:iv:theorem:4:6:6}{第 IV 篇，4.6.6}，\(F\) 可表，当且仅当商 \(F'/F''\)（其中 \(F''=F'\times_F F'\)）在 \((\mathbf{Sch})\) 中存在，并对忠实平坦拟紧态射普遍有效。等价地，商 \(F'/\Gamma\) 必须存在且对这类态射普遍有效。由于 \(F'\) 在 \(S\) 上仿射，\SGARefLink{thm:sga3-V-4.1}{第 V 篇，4.1} 表明此条件满足。这就证明了 \textit{\SGARefLink{sga3:xi:proposition:3:12:item:a}{(a)}} 中 \(F\) 的可表性。

最后，假设 \(G\) 在 \(S\) 上有限呈示。附加断言立即由上述证明得到：由 \textit{\SGARefLink{sga3:xi:proposition:3:12:item:b}{(b)}}，此时 \(F'\) 在 \(S\) 上局部有限呈示，因而可以应用第五讲。\footnote{编者注：原文将最后这处对第五讲的引用标为尚需核查或进一步明确。}


